\clearpage
\fancyhead[L]{\small VOYAGE DE DENTRECASTEAUX.}
\hypertarget{native-519-lower}{}\label{native-519-lower}
\begin{center}\small Göttingische gelehrte Anzeigen. Stück 167. S. 1657..1663. 21. Oktober 1809.\\[3pt]\rule{.72\linewidth}{.3pt}\end{center}

\emph{Voyage de} \textsc{Dentrecasteaux}\emph{, envoyé à la recherche de la Pérouse. Publié par ordre de Sa Majesté l'Empereur et Roi.} \emph{Rédigé par} M. \textsc{de Rossel}\emph{, ancien capitaine de vaisseau.} \emph{Tome second.} \emph{A Paris} 1808. 692 Seiten im Quartformat. Paris. Der erste Teil dieses Werks, der die Reise selbst schildert, wurde in diesen Blättern bereits besprochen (oben S. 1065). Dieser zweite Band ist ganz den astronomisch-nautischen Beobachtungen gewidmet und gliedert sich in zwei Hauptteile. Der erste beschreibt die Instrumente und untersucht allgemein, welche Genauigkeit sich mit ihnen erreichen lässt und unter welchen Umständen man die Beobachtungen am besten anstellt. Der zweite stellt die während der Reise auf der Fregatte La Recherche und zum Teil auch auf der anderen Fregatte, L'Esperance, gemachten astronomischen Beobachtungen in größter Ausführlichkeit und in Tabellen geordnet dar. Die für die Sicherheit der Schifffahrt so wichtigen astronomischen Beobachtungen und die Aufgabe, daraus Breite und Länge zu bestimmen, fallen meist Personen zu, die keine eigentlich mathematische Ausbildung erhalten und sich gewöhnlich kaum mehr als eine mechanische Fertigkeit im Beobachten und Rechnen angeeignet haben. Das ist der Hintergrund für die Abhandlung, die den ersten Teil ausmacht. Sie stammt von einem so sachkundigen und erfahrenen Beobachter wie Kapitän \textsc{Rossel} und ist gewiss eine sehr wertvolle Arbeit. Angehenden Seeleuten ist ihr Studium sehr zu empfehlen, damit sie den Wert besser beurteilen lernen, den sie 'den Beobachtungen beimessen können. Auch Liebhaber, die an Land beobachten, werden sie mit Gewinn lesen. Zu wünschen wäre nur, dass diese Abhandlung nicht Teil eines so teuren Werks wäre, mit dem sie eigentlich in keiner unmittelbaren Verbindung steht, oder dass sie wenigstens durch einen Sonderdruck mehr Lesern zugänglich gemacht würde. Wir hätten uns auch kleinere Unterabschnitte gewünscht, die die Übersicht erleichtern. Es gibt zwar Abschnitte, aber Abschnitte von 119 Seiten Länge. Von den beiden Kapiteln der Abhandlung handelt das erste von den Instrumenten, mit denen die beiden Fregatten ausgerüstet waren. Es erläutert zugleich die besonderen Vorzüge jedes Instruments. Jede Fregatte führte außer mehreren \textsc{Bordai}'schen Reflexionskreisen einen astronomischen Repetitionskreis für den Gebrauch an Land, einen Seechronometer, eine astronomische Uhr, einen Zähler, ein achromatisches Fernrohr, einen Inklinations\SourceLeaf{526}kompass neuer Bauart und einen Azimutkompass mit. Mit den Inklinationskompassen wurden an mehreren Orten auch Versuche über die Schwingungsdauer der Nadel im magnetischen Meridian angestellt. \textsc{Rossel} fügt eine Tafel bei, mit der sich die Schwingungen auf unendlich kleine reduzieren lassen. Der ganze Apparat befand sich in einem mit Glasscheiben verschlossenen Gehäuse, das bei diesen Versuchen jedoch geöffnet werden musste. \textsc{Rossel} meint daher, die Bewegung der Luft könne die Versuche bisweilen etwas unsicher gemacht haben. Schließlich wird noch die Einrichtung eines astronomischen Zeltes beschrieben, wie es jedes Schiff für die Beobachtungen an Land mitführte. — Das zweite Kapitel erläutert die Art und die Grenzen der Fehler, die von den Instrumenten, den Beobachtungen selbst oder den Tafeln herrühren, und die Mittel, sie zu korrigieren oder wenigstens möglichst zu verringern. Die verschiedenen Arten von Beobachtungen werden hier in 9 Abschnitten behandelt. §. I. Über die Zeitbestimmung aus einzelnen Höhen. Der Verfasser weist hier unter anderem auf den Einfluss der Fehler bei der Bestimmung der Horizontsenkung hin. Er empfiehlt daher, stets so hoch über der Meeresoberfläche zu beobachten, wie es praktisch möglich ist. Dieselbe Bemerkung gilt offenbar für alle auf See beobachteten Höhenwinkel. §. II. Breitenbestimmungen aus Meridianhöhen. Sonnenhöhen nahe dem Zenit sind nach \textsc{Rossel} nur auf 2..3 Minuten genau. Der Grund, aus dem er die Messung unter solchen Umständen für weniger zuverlässig hält, leuchtet uns nicht recht ein. Beim Wiegen des Reflexionskreises, so sagt er, löse sich das Bild der Sonne nur sehr langsam vom Seehorizont, sodass sich der eigentliche Vertikalkreis der Berührung nicht genau bestimmen lasse. Sternhöhen auf See zu messen gelang nie. Auch unter den günstigsten Umständen ließen sich Fehler von bis zu 5 Minuten nicht vermeiden. Das liegt vielleicht zum Teil an den unvollkommenen optischen Teilen der französischen Instrumente, deren Fernrohre eine viel zu kleine Öffnung haben. Bei Breitenbestimmungen aus Mondhöhen rechnet \textsc{Rossel} tagsüber mit einer Genauigkeit von 3 und nachts von 4 Minuten. Darin ist aber 1 Minute für einen Fehler der den Ephemeriden entnommenen Monddeklination eingerechnet. Dieser Fehler ließe sich leicht vermeiden, wenn man sich die Mühe macht, die Deklination aus der Länge und Breite zu berechnen, die in den Ephemeriden gewöhnlich auf Sekunden genau angegeben werden. §. III. Breitenbestimmung aus zwei Sonnenhöhen außerhalb des Meridians. Mit Recht wird den indirekten Lösungsmethoden der Vorzug gegeben, auch wenn die direkte Lösung nicht unbedingt so langwierige und mühsame Rechnungen und eine so verwickelte Unterscheidung vieler Fälle erfordert, wie \textsc{Rossel} glaubt. Über die Lösung selbst findet man hier zwar nichts Neues. Sehr ausführlich wird aber erörtert, wie sich Fehler in den Beobachtungen, in der Schätzung der Bewegung des Schiffes zwischen den Beobachtungen und im Gang der Uhr auswirken. Das Endergebnis ist, dass man mit diesem Verfahren die Breite stets auf 4 bis 5 Minuten genau bestimmen könne, selbst im schlimmsten Fall, wenn alle Fehler in derselben Richtung zusammenwirken, sofern man bei der Auswahl der Beobachtungen die nötige Vorsicht walten lässt. Was über die größere oder geringere Genauigkeit der verschiedenen indirekten Rechenverfahren gesagt wird, dürfte für die Fälle, die in der Praxis vorkommen, ziemlich überflüssig sein. §. IV. Beobachtung des Azimuts der Sonne zur Bestimmung der Abweichung der Magnetnadel. \textsc{Rossel} weist darauf hin, dass die an Bord vor der Küste gemessene Abweichung bisweilen erheblich von der in geringer Entfernung an Land gefundenen abweicht. Die Beobachtungen bei Teneriffa, die wir weiter unten anführen, sind dafür ein bemerkenswertes Beispiel. §. V. Längenbestimmungen aus gemessenen Abständen des Mondes von der Sonne und von Fixsternen. Hier findet man manche wertvolle praktische Bemerkung, die als Ergebnis langer Erfahrung umso mehr Gewicht hat. Der Verfasser sagt gewiss mit Recht: „Alles, was eine so heikle Beobachtung erleichtern kann, ist wertvoll. Wir kennen \SourceLeaf{527}keine, die mehr natürliche Begabung erfordert. Daher sind wir fest davon überzeugt, dass das Geschick des Beobachters noch viel beitragen wird, wenn die Gelehrten alles getan haben, was ihnen ihr Genie zur Vervollkommnung dieser Methode eingeben kann.“ — Bei der Berechnung der Mondparallaxe zur Reduktion der gemessenen Abstände berücksichtigt \textsc{Rossel} die sphäroidische Gestalt der Erde nur insofern, als die Horizontalparallaxe auf der Breite des Beobachtungsortes kleiner ist als die Äquatorialparallaxe. Die Neigung der Vertikalen gegen den zum Erdmittelpunkt gerichteten Radius vernachlässigt er. Das hat bei Beobachtungen, wie sie sich mit gewöhnlichen französischen Spiegelkreisen anstellen lassen, freilich wenig zu bedeuten. Den Einfluss dieses letzten Umstands hat \textsc{Rossel} jedoch unserer Meinung nach falsch beurteilt. Er berücksichtigt nur, dass die Parallaxe für eine falsche \emph{Höhe} berechnet wird, und vergisst, dass sich dadurch auch die \emph{Richtung} der Parallaxe ändert. Daher halten wir seinen Schluss für irrig, das Ergebnis werde durch dieses nicht ganz richtige Rechenverfahren am stärksten verfälscht, wenn der Mond nahe der Kulmination steht. In der Regel wird der Fehler gerade dann am größten, wenn sich der Mond nahe dem ersten Vertikal befindet. §. VI. Längenbestimmungen durch Seechronometer. §. VII. Wie die durch Mondabstände bestimmten Längen mit den durch Seechronometer erhaltenen Längenunterschieden zu verbinden sind. §. VIII. Längenbestimmungen aus Sternbedeckungen und Finsternissen der Jupitermonde. §. IX. Bestimmung der Azimute irdischer Gegenstände (relevemens).

Der zweite Teil des Werks enthält das Beobachtungsjournal selbst. Ihm gehen eine Tafel der Fehler der \textsc{Mason}'schen Mondtafeln nach gleichzeitigen Beobachtungen in Greenwich und die daraus abgeleiteten Korrekturen der während der Reise bestimmten Längen voraus. Einzelne der für die Geographie der Südsee sehr wichtigen Beobachtungsergebnisse hier herauszugreifen, würde dem Plan dieser \emph{Göttingischen gelehrten Anzeigen} zuwiderlaufen. Es wäre umso weniger nötig, als sich ein sehr vollständiger Auszug der geographischen Bestimmungen bereits in der \emph{Monatlichen Correspondenz} findet. Manchem Leser ist es aber vielleicht willkommen, wenn wir hier die wenigen Bestimmungen mit der Magnetnadel mitteilen, die in jener Zeitschrift übergangen wurden. In \emph{Brest} fand man am 20. September 1791 für die Neigung der Nadel nach Norden $71^\circ\,30'$ und für die Dauer einer unendlich kleinen Schwingung $2^{\mathrm s}02$. In \emph{Santa Cruz} auf Teneriffa betrug die Neigung nach Norden am 21. Oktober 1791 $62^\circ\,25'$ und die mittlere Schwingungsdauer $2^{\mathrm s}081$. Die Deklination der Nadel betrug an Bord der Recherche $18^\circ\,7'$ N. W., beim Observatorium $21^\circ\,33'$ und auf der Mole $23^\circ\,43'$. In \emph{Port du Nord} auf Van-Diemens-Land betrug am 11. Mai 1792 die Neigung $70^\circ\,50'$ nach Süden, die mittlere Schwingungsdauer $1^{\mathrm s}869$ und die Abweichung $8^\circ\,1'\,12''$ N. O. Auf der Insel \emph{Amboina} betrug die Neigung am 9. Oktober 1792 $20^\circ\,37'$ nach Süden und die Schwingungsdauer $2^{\mathrm s}403$. In \emph{Port au Sud} auf Van-Diemens-Land betrug die Abweichung am 7. Februar 1793 mit der einen Nadel $2^\circ\,33'\,47''$ N. O. und mit der anderen $3^\circ\,10'\,24''$. Die Neigungen betrugen $72^\circ\,22'$ und $70^\circ\,48'$ nach Süden, die Schwingungsdauern $1^{\mathrm s}8498$ und $1^{\mathrm s}817$. In \emph{Sourabaya} auf der Insel Java betrug am 9. Mai 1794 die Neigung $25^\circ\,40'$ nach Süden und die Schwingungsdauer mit der einen Nadel $2^{\mathrm s}429$, mit der anderen $2^{\mathrm s}485$. Aus den verschiedenen Beobachtungen der Schwingungsdauer scheint also hervorzugehen, dass die Stärke der magnetischen Kraft am Äquator beträchtlich geringer ist als in höheren Breiten.
